\section{Rekomendasi}
\label{sec:rekomendasi}

\subsection{Langkah Konkret}

\begin{enumerate}
  \item \textbf{Prioritaskan IPP tua dengan dampak kesehatan tinggi.} Pensiun dini paling layak secara keuangan untuk PLTU IPP yang pembayaran kapasitasnya masih berjalan. Prioritaskan unit yang juga memiliki biaya kesehatan per kWh tinggi, seperti Cirebon-1.
  \item \textbf{Pastikan kompensasi didanai dari luar neraca PLN.} H2 hanya berlaku bila sebagian besar kompensasi berasal dari ETM/JETP atau dana iklim. Tanpa itu, pensiun dini memindahkan beban ke tarif dan APBN.
  \item \textbf{Hitung pembangkit pengganti sejak awal.} Jendela oversupply Jawa-Bali praktis menutup pada 2027-2028. Setiap usulan pensiun dini setelahnya perlu memasukkan biaya pengganti bersih, misalnya PLTS skala besar dengan baterai.
  \item \textbf{Dorong PLTS atap lewat jalur yang tidak memukul pendapatan PLN.} Jalur yang paling seimbang adalah gedung pemerintah, karena tagihan yang dihemat kembali ke APBN. Mengembalikan net metering secara luas membutuhkan mekanisme kompensasi pendapatan PLN, misalnya dari dana iklim yang dihitung dari nilai kesehatan.
  \item \textbf{Gunakan alat terbuka dalam kajian Permen ESDM 10/2025.} Kajian enam bulan PLN dan evaluasi tim kerja gabungan sebaiknya memakai asumsi yang dapat diperiksa publik, misalnya lewat PowerShift.
\end{enumerate}

\subsection{Peran Aktor}

\begin{table}[H]
\centering
\caption{Pembagian peran aktor}
\label{tab:aktor}
\small
\begin{tabularx}{\textwidth}{p{0.24\textwidth} X}
\toprule
Aktor & Tindakan \\
\midrule
Kementerian ESDM & Menetapkan daftar kandidat pensiun dini berbasis kriteria Pasal 11 dan biaya kesehatan, memakai alat terbuka dalam evaluasi tim kerja gabungan \\
PLN & Menyediakan data pembayaran kapasitas dan akhir PPA untuk kajian, menghitung kebutuhan pembangkit pengganti \\
Kementerian Keuangan & Menyusun skema kompensasi dari dana luar dan menghitung penghematan subsidi serta tagihan pemerintah \\
Mitra ETM/JETP dan ADB & Mengalihkan dana Cirebon-1 ke paket IPP tua dengan dampak kesehatan tinggi \\
Pemerintah daerah & Memasang PLTS atap di gedung pemerintah dan menyiapkan program alih keterampilan pekerja PLTU \\
Akademisi dan masyarakat sipil & Memelihara data dan model PowerShift, menyosialisasikan hasil kepada publik \\
\bottomrule
\end{tabularx}
\end{table}

Secara keseluruhan, transisi PLTS atap dan pensiun dini PLTU macet karena manfaatnya dirasakan masyarakat (kesehatan, tagihan) sementara biayanya jatuh ke PLN dan APBN. Paket yang memilih PLTU dengan tepat, mendanai kompensasi dari luar, dan memasang PLTS atap lewat jalur pemerintah dapat menyeimbangkan kedua sisi itu.
